\documentclass[10pt]{article}
\usepackage[T1]{fontenc}
\usepackage[utf8]{inputenc}
\usepackage[french]{babel}
\usepackage{amsmath,amssymb,mathrsfs,mathpazo}
\usepackage{geometry}
\geometry{a4paper,margin=2cm}
\usepackage{graphicx,array,tabularx,tabularray,colortbl,multirow,multicol,booktabs,enumitem,pifont,ulem,fancybox,diagbox,cancel}
\usepackage{pstricks,pstricks-add,pst-plot,pst-tree,pst-node,pst-text,pst-func,pst-math,pst-bspline,pst-3dplot,pst-eucl,pst-solides3d}
\let\pscalculate\undefined
\usepackage{pst-bezier}
\usepackage{tikz,pgfplots,tkz-tab}
\usetikzlibrary{arrows,arrows.meta,patterns,calc,positioning,trees,decorations.markings,shapes.geometric,shapes.misc,shapes.arrows,intersections,decorations.pathmorphing,decorations.pathreplacing,backgrounds,shadows,babel}
\pgfplotsset{compat=1.18}
\usepgfplotslibrary{fillbetween}
\usepackage[np]{numprint}
\usepackage{listings,makecell,pas-tableur,twoopt}
\let\barre\undefined
\usepackage[pstricks]{tablvar}
\definecolor{darkgreen}{RGB}{0,100,0}
\providecommand{\red}{\color{red}}
\providecommand{\blue}{\color{blue}}
\providecommand{\green}{\color{green}}
\providecommand{\black}{\color{black}}
\providecommand{\white}{\color{white}}
\providecommand{\gray}{\color{gray}}
\providecommand{\cyan}{\color{cyan}}
\providecommand{\magenta}{\color{magenta}}
\providecommand{\yellow}{\color{yellow}}
\providecommand{\violet}{\color{violet}}
\providecommand{\orange}{\color{orange}}
\providecommand{\euro}{\texteuro}
\providecommand{\EUR}{\texteuro}
\providecommand{\ieme}{\textsuperscript{e}}
\providecommand{\degres}{\ensuremath{^\circ}}
\providecommand{\degree}{\ensuremath{^\circ}}
\providecommand{\vv}[1]{\overrightarrow{#1}}
\providecommand{\Coord}[1]{#1}
\providecommand{\decosix}{}
\providecommand{\starredbullet}{\ensuremath{\bullet}}
\pagestyle{empty}
\setlength{\parindent}{0pt}
\usepackage{adjustbox}\DeclareUnicodeCharacter{2236}{:}\DeclareUnicodeCharacter{2190}{\ensuremath{\leftarrow}}
\providecommand{\R}{\mathbb{R}}
\providecommand{\N}{\mathbb{N}}
\providecommand{\Z}{\mathbb{Z}}
\providecommand{\Q}{\mathbb{Q}}
\providecommand{\D}{\mathbb{D}}
\providecommand{\C}{\mathbb{C}}
\providecommand{\e}{{\mathrm{e}}}
\providecommand{\vect}[1]{\overrightarrow{#1}}
\providecommand{\vectt}[1]{\overrightarrow{\mathrm{#1}}}
\providecommand{\vvt}[1]{\overrightarrow{\mathrm{#1}}}
\providecommand{\np}[2][]{{#2~#1}}
\providecommand{\vv}[1]{{\overrightarrow{#1}}}
\renewcommand{\d}{{\mathrm{d}}}
\providecommand{\degres}{{^\circ}}
\providecommand{\degre}{{^\circ}}
\providecommand{\blue}{\color{blue}}
\providecommand{\red}{\color{red}}
\providecommand{\enskip}{\quad}
\providecommand{\llbracket}{\lbrack\!\lbrack}
\providecommand{\rrbracket}{\rbrack\!\rbrack}
\providecommand{\euro}{{\text{€}}}
\providecommand{\no}{{\text{n°}}}
\providecommand{\interval}[2]{{[#1;#2]}}
\providecommand{\intervFO}[2]{{[#1;#2[}}
\providecommand{\dotfill}{\ldots}

\providecommand{\ds}{\displaystyle}
\providecommand{\textding}[1]{#1}
\providecommand{\Oij}{\(\mathrm{O};\vec{\imath},\vec{\jmath}\)}
\providecommand{\Oijk}{\(\mathrm{O};\vec{\imath},\vec{\jmath},\vec{k}\)}
\providecommand{\Ouv}{\(\mathrm{O};\vec{u},\vec{v}\)}
\providecommand{\intervFF}[2]{\ensuremath{[#1;#2]}}
\renewcommand{\euro}{\ifmmode\text{\texteuro}\else\texteuro\fi}
\DeclareUnicodeCharacter{00B0}{\ifmmode^{\circ}\else\textdegree\fi}
\DeclareUnicodeCharacter{20AC}{\ifmmode\text{\texteuro}\else\texteuro\fi}
% Énoncés : Éducation nationale. Transcription et corrigés : APMEP, Denis Vergès et contributeurs. Découpage et classement : CoopMaths, Sébastien Lozano, Éric Elter et contributeurs. Adaptation de la présentation : Maths Explorer.
% CC BY-SA : https://coopmaths.fr/www/about/
% https://coopmaths.fr/alea/static/e3c/2020/tex/e3c_2020_00_sujet31_2_cor.tex
\begin{document}
\noindent\textit{APMEP / CoopMaths. CC BY-SA. Présentation adaptée par Maths Explorer.}\par\medskip
\(f(x) = (5 - 2x)\e^x\).

\subsection*{1.}

\begin{itemize}
    \item \(A\) d'abscisse nulle appartient à \(C\), son ordonnée est donc \(f(0) = 5\e^0 = 5\). Donc \(A(0\,;\,5)\).
    \item \(B\) d'ordonnée nulle appartient à \(C\), son abscisse est donc telle que :
    \[
    f(x) = 0 \quad \Longleftrightarrow \quad (5 - 2x)\e^x = 0 \quad \Longleftrightarrow \quad 5 - 2x = 0,
    \]
    puisqu'on sait que, quel que soit \(x\),  \(\e^x > 0\).
    
    Donc : \(5 = 2x \Longleftrightarrow x = \dfrac{5}{2}\) et \(B\left(\dfrac{5}{2}\,;\,0\right)\).
\end{itemize}

\subsection*{2.}

\(f\), produit de fonctions dérivables sur \(\mathbb{R}\), est dérivable sur \(\mathbb{R}\) et sur cet intervalle :
\[
f'(x) = -2 \e^x + (5 - 2x)\e^x = \e^x(-2 + 5 - 2x) = (3 - 2x)\e^x.
\]

\subsection*{3.}

\(f'(x)\) est donc du signe de \(3 - 2x\) :
\begin{align*}
&3 - 2x > 0 \\
\iff&  3 > 2x \\
\iff & \dfrac{3}{2} > x
\end{align*}
donc la fonction \(f\) est croissante sur \(\left]-\infty\,;\,\dfrac{3}{2}\right[\) ;
	\begin{center}
	\begin{tikzpicture}[double distance=2pt]
		\tkzTabInit{$x$/1,f'(x)/1,$f(x)$/2}{$0$,$\dfrac{3}{2}$, $10$}
		\tkzTabLine{,+,z, -}
		\tkzTabVar{+/$0$,-/ $6\e^{-1}$ /, +/}
	\end{tikzpicture}
\end{center} 

\subsection*{4.}

D'après le résultat précédent \(D\left(1,5\,;\,2 e^{1,5}\right)\).

\subsection*{5.}

Une équation de la tangente \(T\) à la courbe \(C\) au point \(A\) d'abscisse \(0\) est :

\[
M(x\,;\,y) \in \mathcal T \quad \Longleftrightarrow \quad y - f(0) = f'(0)(x - 0).
\]

Avec \(f(0) = 5\) et \(f'(0) = 3\), on obtient :
\[
M(x\,;\,y) \in \mathcal T \quad \iff \quad y - 5 = 3x \quad \iff \quad y = 3x + 5.
\]

Or $D\left(1,5\,;\,2 \e^{1,5}\right) \in \mathcal T \iff 2\e^{1,5} = 4,5 + 5 $\\
Comme $ 2 \e^{1,5} \neq 9,5$ , alors $ D \notin \mathcal T$.

\end{document}
